\documentclass[../cours_tle.tex]{subfiles}
\begin{document}
\setcounter{chapter}{12}
\chapter{Continuité, valeurs intermédiaires et équations}
\label{ch:13}
\begin{objectif}
Spécialité. Prouver qu'une solution existe puis établir son unicité. Prérequis : chapitres 9 à 12.
\end{objectif}
\section{Continuité en un point}
\begin{definition}
Une fonction $f$ définie en $a$ est continue en $a$ si $f(x)\to f(a)$ quand $x\to a$ dans son domaine. Elle est continue sur un intervalle si elle l'est en chacun de ses points, avec une limite unilatérale aux extrémités.
\end{definition}
Les sommes, produits, composées de fonctions continues sont continus là où ils sont définis ; les quotients le sont lorsque le dénominateur ne s'annule pas, par les règles de limites.
Un saut est une discontinuité. Une fonction peut aussi ne pas être définie en un point mais y posséder une limite finie : définir la valeur manquante égale à cette limite réalise un prolongement par continuité.
\begin{theoreme}
Si $f$ est continue en $\ell$ et $u_n\to\ell$, alors $f(u_n)\to f(\ell)$, pourvu que les $u_n$ soient dans son domaine.
\end{theoreme}
\begin{demonstration}
Fixer $\varepsilon>0$. La continuité fournit $\delta>0$ tel que $|x-\ell|<\delta$ implique $|f(x)-f(\ell)|<\varepsilon$. La convergence de $u_n$ assure la première inégalité à partir d'un rang ; on obtient alors la seconde.
\end{demonstration}
\section{Le théorème des valeurs intermédiaires}
\begin{theoreme}
Si $f$ est continue sur $[a,b]$, tout nombre $k$ entre $f(a)$ et $f(b)$ est l'image d'au moins un $c\in[a,b]$.
Si $f$ est en outre strictement monotone, cet antécédent est unique.
\end{theoreme}
\begin{demonstration}[Par dichotomie, approfondissement]
Considérons $g=f-k$, avec $g(a)\le0\le g(b)$ ; les autres sens se traitent en remplaçant $g$ par $-g$. Si une borne est racine, terminer. Sinon couper l'intervalle au milieu et conserver une moitié dont les bornes ont des signes opposés (ou terminer si le milieu est racine).
Les bornes $a_n$ croissent, les $b_n$ décroissent et $b_n-a_n=(b-a)/2^n\to0$. Elles tendent donc vers un même $c$. Par continuité, $g(a_n)\to g(c)$ et $g(b_n)\to g(c)$ ; leurs signes imposent $g(c)=0$.
Une fonction strictement monotone est injective, ce qui donne l'unicité.
\end{demonstration}
Sans continuité, une fonction qui vaut $-1$ à gauche de $0$ et $1$ à partir de $0$ n'atteint pas $0$. Sans monotonie stricte, une valeur peut être atteinte plusieurs fois.
\section{Résoudre : existence, unicité, approximation}
Pour $f(x)=x^3+x-1$ sur $[0,1]$, $f(0)=-1$, $f(1)=1$. La continuité donne une racine ; $f'(x)=3x^2+1>0$ donne son unicité. Puis $f(0{,}68)<0<f(0{,}69)$ donne un encadrement.
Après $n$ dichotomies, choisir le milieu garantit une erreur au plus $(b-a)/2^{n+1}$.
On ne remplace pas la preuve par la valeur affichée par une calculatrice.
\section{Complément : réciproque d'une fonction monotone}
Une fonction continue strictement croissante sur un intervalle est une bijection sur son image, par le théorème précédent. Sa réciproque est continue : pour un point intérieur $x_0$, choisir $x_-<x_0<x_+$ aussi proches que souhaité ;
$f(x_-)<f(x_0)<f(x_+)$ encadre les images, et l'ordre inverse encadre les antécédents de tout $y$ assez proche de $f(x_0)$. Aux extrémités, utiliser un encadrement unilatéral.
Cette preuve servira au logarithme, sans admettre sa continuité séparément.
\section{Complément : une équation fonctionnelle}
Si $f:\R\to\R$ est continue et $f(x+y)=f(x)+f(y)$, alors $f(x)=cx$, où $c=f(1)$.
En effet $f(0)=0$, $f(-x)=-f(x)$, puis par addition $f(n)=nc$ et $f(p/q)=pc/q$ pour $q>0$. Tout réel est limite de rationnels ; la continuité étend la formule à $\R$. La continuité est une hypothèse essentielle à cette extension.
\begin{finchapitre}
\section{Utilité et suite du parcours}
Le théorème des valeurs intermédiaires est un théorème d'existence. La distinction entre existence, unicité et calcul approché se retrouve dans les équations différentielles et les problèmes de recherche.
\begin{retenir}
Pour une équation, annoncer l'intervalle et vérifier continuité, encadrement de la valeur cible et monotonie pour l'unicité.
\end{retenir}
\end{finchapitre}
\end{document}
